% Options for packages loaded elsewhere
\PassOptionsToPackage{unicode}{hyperref}
\PassOptionsToPackage{hyphens}{url}
\documentclass[
]{article}
\usepackage{xcolor}
\usepackage{amsmath,amssymb}
\setcounter{secnumdepth}{-\maxdimen} % remove section numbering
\usepackage{iftex}
\ifPDFTeX
  \usepackage[T1]{fontenc}
  \usepackage[utf8]{inputenc}
  \usepackage{textcomp} % provide euro and other symbols
\else % if luatex or xetex
  \usepackage{unicode-math} % this also loads fontspec
  \defaultfontfeatures{Scale=MatchLowercase}
  \defaultfontfeatures[\rmfamily]{Ligatures=TeX,Scale=1}
\fi
\usepackage{lmodern}
\ifPDFTeX\else
  % xetex/luatex font selection
\fi
% Use upquote if available, for straight quotes in verbatim environments
\IfFileExists{upquote.sty}{\usepackage{upquote}}{}
\IfFileExists{microtype.sty}{% use microtype if available
  \usepackage[]{microtype}
  \UseMicrotypeSet[protrusion]{basicmath} % disable protrusion for tt fonts
}{}
\makeatletter
\@ifundefined{KOMAClassName}{% if non-KOMA class
  \IfFileExists{parskip.sty}{%
    \usepackage{parskip}
  }{% else
    \setlength{\parindent}{0pt}
    \setlength{\parskip}{6pt plus 2pt minus 1pt}}
}{% if KOMA class
  \KOMAoptions{parskip=half}}
\makeatother
\usepackage{color}
\usepackage{fancyvrb}
\newcommand{\VerbBar}{|}
\newcommand{\VERB}{\Verb[commandchars=\\\{\}]}
\DefineVerbatimEnvironment{Highlighting}{Verbatim}{commandchars=\\\{\}}
% Add ',fontsize=\small' for more characters per line
\newenvironment{Shaded}{}{}
\newcommand{\AlertTok}[1]{\textcolor[rgb]{1.00,0.00,0.00}{\textbf{#1}}}
\newcommand{\AnnotationTok}[1]{\textcolor[rgb]{0.38,0.63,0.69}{\textbf{\textit{#1}}}}
\newcommand{\AttributeTok}[1]{\textcolor[rgb]{0.49,0.56,0.16}{#1}}
\newcommand{\BaseNTok}[1]{\textcolor[rgb]{0.25,0.63,0.44}{#1}}
\newcommand{\BuiltInTok}[1]{\textcolor[rgb]{0.00,0.50,0.00}{#1}}
\newcommand{\CharTok}[1]{\textcolor[rgb]{0.25,0.44,0.63}{#1}}
\newcommand{\CommentTok}[1]{\textcolor[rgb]{0.38,0.63,0.69}{\textit{#1}}}
\newcommand{\CommentVarTok}[1]{\textcolor[rgb]{0.38,0.63,0.69}{\textbf{\textit{#1}}}}
\newcommand{\ConstantTok}[1]{\textcolor[rgb]{0.53,0.00,0.00}{#1}}
\newcommand{\ControlFlowTok}[1]{\textcolor[rgb]{0.00,0.44,0.13}{\textbf{#1}}}
\newcommand{\DataTypeTok}[1]{\textcolor[rgb]{0.56,0.13,0.00}{#1}}
\newcommand{\DecValTok}[1]{\textcolor[rgb]{0.25,0.63,0.44}{#1}}
\newcommand{\DocumentationTok}[1]{\textcolor[rgb]{0.73,0.13,0.13}{\textit{#1}}}
\newcommand{\ErrorTok}[1]{\textcolor[rgb]{1.00,0.00,0.00}{\textbf{#1}}}
\newcommand{\ExtensionTok}[1]{#1}
\newcommand{\FloatTok}[1]{\textcolor[rgb]{0.25,0.63,0.44}{#1}}
\newcommand{\FunctionTok}[1]{\textcolor[rgb]{0.02,0.16,0.49}{#1}}
\newcommand{\ImportTok}[1]{\textcolor[rgb]{0.00,0.50,0.00}{\textbf{#1}}}
\newcommand{\InformationTok}[1]{\textcolor[rgb]{0.38,0.63,0.69}{\textbf{\textit{#1}}}}
\newcommand{\KeywordTok}[1]{\textcolor[rgb]{0.00,0.44,0.13}{\textbf{#1}}}
\newcommand{\NormalTok}[1]{#1}
\newcommand{\OperatorTok}[1]{\textcolor[rgb]{0.40,0.40,0.40}{#1}}
\newcommand{\OtherTok}[1]{\textcolor[rgb]{0.00,0.44,0.13}{#1}}
\newcommand{\PreprocessorTok}[1]{\textcolor[rgb]{0.74,0.48,0.00}{#1}}
\newcommand{\RegionMarkerTok}[1]{#1}
\newcommand{\SpecialCharTok}[1]{\textcolor[rgb]{0.25,0.44,0.63}{#1}}
\newcommand{\SpecialStringTok}[1]{\textcolor[rgb]{0.73,0.40,0.53}{#1}}
\newcommand{\StringTok}[1]{\textcolor[rgb]{0.25,0.44,0.63}{#1}}
\newcommand{\VariableTok}[1]{\textcolor[rgb]{0.10,0.09,0.49}{#1}}
\newcommand{\VerbatimStringTok}[1]{\textcolor[rgb]{0.25,0.44,0.63}{#1}}
\newcommand{\WarningTok}[1]{\textcolor[rgb]{0.38,0.63,0.69}{\textbf{\textit{#1}}}}
\usepackage{longtable,booktabs,array}
\newcounter{none} % for unnumbered tables
\usepackage{calc} % for calculating minipage widths
% Correct order of tables after \paragraph or \subparagraph
\usepackage{etoolbox}
\makeatletter
\patchcmd\longtable{\par}{\if@noskipsec\mbox{}\fi\par}{}{}
\makeatother
% Allow footnotes in longtable head/foot
\IfFileExists{footnotehyper.sty}{\usepackage{footnotehyper}}{\usepackage{footnote}}
\makesavenoteenv{longtable}
\setlength{\emergencystretch}{3em} % prevent overfull lines
\providecommand{\tightlist}{%
  \setlength{\itemsep}{0pt}\setlength{\parskip}{0pt}}
\usepackage{bookmark}
\IfFileExists{xurl.sty}{\usepackage{xurl}}{} % add URL line breaks if available
\urlstyle{same}
\hypersetup{
  hidelinks,
  pdfcreator={LaTeX via pandoc}}

\author{}
\date{}

\begin{document}

\section{Paper 2: Radius Sweep of Fully-Inscribed Unit-Cell Counts on a
4-Dimensional
Lattice}\label{paper-2-radius-sweep-of-fully-inscribed-unit-cell-counts-on-a-4-dimensional-lattice}

\subsection{An Enumeration Table from Radius 0.5 to 10.0 and a
Reproducible
Formulation}\label{an-enumeration-table-from-radius-0.5-to-10.0-and-a-reproducible-formulation}

\textbf{Author}: Noriaki Kihara\\
\textbf{Affiliation}: WF System Co., Ltd.\\
\textbf{ORCID}: 0009-0004-6753-4020\\
\textbf{Version}: v0.2\\
\textbf{Date}: June 2026\\
\textbf{DOI (this version)}: 10.5281/zenodo.20607574\\
\textbf{Concept DOI}: 10.5281/zenodo.20588038\\
\textbf{License}: CC BY 4.0

\begin{center}\rule{0.5\linewidth}{0.5pt}\end{center}

\subsection{Abstract}\label{abstract}

In this paper we count the number of 4-dimensional unit cells of side 1,
placed on a 4-dimensional integer lattice, that are fully inscribed in a
4-dimensional hyperball of radius \(R\). The object is a purely
geometric / integer-lattice counting problem, and we do not treat any
correspondence with physical constants.

Taking the center of a unit cell to be

\[
k=(k_1,k_2,k_3,k_4)\in\mathbb{Z}^4 ,
\]

with half-width \(1/2\) in each direction, the condition that this unit
cell be fully inscribed in the hyperball of radius \(R\) is

\[
\sum_{i=1}^{4}
\left(|k_i|+\frac12\right)^2
\le R^2 .
\]

We define the number of lattice points satisfying this condition as

\[
N_0(R)
=
\#\left\{
k\in\mathbb{Z}^4
\mid
\sum_{i=1}^{4}
\left(|k_i|+\frac12\right)^2
\le R^2
\right\} .
\]

In this paper we compute \(N_0(R)\) for radii
\(R=0.5,1.0,1.5,\ldots,10.0\), and tabulate the diameter \(2R\) of the
circumscribing ball and the diagonal length \(2\rho(R)\) of the actually
stacked cell set.

\begin{center}\rule{0.5\linewidth}{0.5pt}\end{center}

\subsection{1. Purpose}\label{purpose}

In Paper 1 we showed that choosing a 4-dimensional lattice in the
wavelength space or the frequency space turns the sum-of-squares
condition into a unit-cell counting problem.

In this paper we actually carry out that counting.

What we treat here is not a physical theory but the following purely
geometric problem.

\begin{quote}
Among the unit 4-cells of side 1 on the 4-dimensional integer lattice,
how many are fully inscribed in the 4-dimensional hyperball of radius
\(R\)?
\end{quote}

\begin{center}\rule{0.5\linewidth}{0.5pt}\end{center}

\subsection{2. Basic Definitions}\label{basic-definitions}

Let the 4-dimensional integer lattice be

\[
\mathbb{Z}^4 .
\]

To each lattice point

\[
k=(k_1,k_2,k_3,k_4)
\]

associate a 4-dimensional unit cell of side 1.

When the cell center is at \(k\), each point \(x_i\) inside the cell
satisfies

\[
k_i-\frac12\le x_i\le k_i+\frac12 .
\]

The squared distance from the origin to the farthest vertex of this cell
is

\[
r_{\max}(k)^2
=
\sum_{i=1}^{4}
\left(|k_i|+\frac12\right)^2 .
\]

Hence the condition that this cell be fully inscribed in the hyperball
of radius \(R\) is

\[
r_{\max}(k)^2\le R^2 ,
\]

that is,

\[
\sum_{i=1}^{4}
\left(|k_i|+\frac12\right)^2
\le R^2 .
\]

\begin{center}\rule{0.5\linewidth}{0.5pt}\end{center}

\subsection{3. Fully-Inscribed Cell
Count}\label{fully-inscribed-cell-count}

We define the fully-inscribed cell count as

\[
N_0(R)
=
\#\left\{
k\in\mathbb{Z}^4
\mid
\sum_{i=1}^{4}
\left(|k_i|+\frac12\right)^2
\le R^2
\right\} .
\]

Multiplying both sides by 4,

\[
\sum_{i=1}^{4}
(2|k_i|+1)^2
\le
(2R)^2 .
\]

Hence this problem can also be expressed as an integer inequality on a
sum of four positive odd squares.

\begin{center}\rule{0.5\linewidth}{0.5pt}\end{center}

\subsection{4. Stacked Diagonal Length}\label{stacked-diagonal-length}

Let the set of cells inscribed at radius \(R\) be

\[
K_R=
\left\{
k\in\mathbb{Z}^4
\mid
r_{\max}(k)\le R
\right\} .
\]

Among the cells in this set, define the distance from the origin to the
farthest vertex as

\[
\rho(R)
=
\max_{k\in K_R} r_{\max}(k) .
\]

If \(K_R\) is empty, we set

\[
\rho(R)=0 .
\]

The centrally symmetric maximal diagonal length of the stacked cell set
is then defined as

\[
L(R)=2\rho(R) ,
\]

and the diameter of the circumscribing ball is

\[
D(R)=2R .
\]

Hence always

\[
L(R)\le D(R) .
\]

\begin{center}\rule{0.5\linewidth}{0.5pt}\end{center}

\subsection{5. Radius Sweep Results}\label{radius-sweep-results}

\subsubsection{5.1 Definitions of Volume, Filling Ratio, and
Gap}\label{definitions-of-volume-filling-ratio-and-gap}

In addition to the fully-inscribed cell count \(N_0(R)\), we define the
following quantities in order to evaluate the degree of filling relative
to the surrounding hyperball.

The volume of the surrounding 4-dimensional hyperball of radius \(R\) is

\[
V_0(R)=\frac{\pi^2}{2}R^4 .
\]

Since each fully-inscribed unit 4-cell has volume \(1^4=1\), the total
volume of the stacked cells is equal to the fully-inscribed cell count:

\[
V_1(R)=N_0(R)\cdot 1^4=N_0(R) .
\]

The filling ratio is therefore defined as

\[
\mathrm{fill}(R)=\frac{V_1(R)}{V_0(R)}=\frac{2\,N_0(R)}{\pi^2 R^4}.
\]

The filling ratio is always \(\mathrm{fill}(R)<1.0\). Over a finite
radius sweep it may show local increases and decreases, while for large
\(R\) it tends toward \(1.0\).

The unfilled part of the surrounding hyperball volume is defined as the
gap:

\[
\mathrm{gap}(R)=V_0(R)-V_1(R)=\frac{\pi^2}{2}R^4-N_0(R)\ (>0).
\]

\subsubsection{5.2 Sweep Results}\label{sweep-results}

Table 1 shows the results computed from radius \(R=0.5\) to \(10.0\) in
steps of \(0.5\), with \(R=100,\,1000,\,10000\) also included. The
values of \(N_0(R)\), \(2\rho(R)\), \(\mathrm{fill}(R)\), and
\(\mathrm{gap}(R)\) are exact consequences of integer arithmetic, with
the check values \(N_0(1{:}5)=1,9,137,473,1545\). The computational
procedure is given in Appendix B.

\textbf{Table 1. Radius sweep of fully-inscribed unit-cell counts on the
4-dimensional lattice, including filling ratio and gap}

{\def\LTcaptype{none} % do not increment counter
\begin{longtable}[]{@{}
  >{\raggedleft\arraybackslash}p{(\linewidth - 12\tabcolsep) * \real{0.1429}}
  >{\raggedleft\arraybackslash}p{(\linewidth - 12\tabcolsep) * \real{0.1429}}
  >{\raggedleft\arraybackslash}p{(\linewidth - 12\tabcolsep) * \real{0.1429}}
  >{\raggedleft\arraybackslash}p{(\linewidth - 12\tabcolsep) * \real{0.1429}}
  >{\raggedleft\arraybackslash}p{(\linewidth - 12\tabcolsep) * \real{0.1429}}
  >{\raggedleft\arraybackslash}p{(\linewidth - 12\tabcolsep) * \real{0.1429}}
  >{\raggedleft\arraybackslash}p{(\linewidth - 12\tabcolsep) * \real{0.1429}}@{}}
\toprule\noalign{}
\begin{minipage}[b]{\linewidth}\raggedleft
Circumscribing radius R
\end{minipage} & \begin{minipage}[b]{\linewidth}\raggedleft
Diameter 2R
\end{minipage} & \begin{minipage}[b]{\linewidth}\raggedleft
Fully-inscribed cell count N0(R)
\end{minipage} & \begin{minipage}[b]{\linewidth}\raggedleft
Stacked diagonal \(2\rho(R)\)
\end{minipage} & \begin{minipage}[b]{\linewidth}\raggedleft
Radial margin \(R-\rho(R)\)
\end{minipage} & \begin{minipage}[b]{\linewidth}\raggedleft
Filling ratio V1/V0
\end{minipage} & \begin{minipage}[b]{\linewidth}\raggedleft
Gap V0-V1
\end{minipage} \\
\midrule\noalign{}
\endhead
\bottomrule\noalign{}
\endlastfoot
0.5 & 1 & 0 & 0 & 0.5 & 0.00000000 & 0.3 \\
1 & 2 & 1 & 2 & 0 & 0.20264237 & 3.9 \\
1.5 & 3 & 1 & 2 & 0.5 & 0.04002812 & 24.0 \\
2 & 4 & 9 & 3.4641 & 0.267949 & 0.11398633 & 70.0 \\
2.5 & 5 & 33 & 4.47214 & 0.263932 & 0.17119227 & 159.8 \\
3 & 6 & 137 & 6 & 0 & 0.34274079 & 262.7 \\
3.5 & 7 & 233 & 6.63325 & 0.183375 & 0.31464004 & 507.5 \\
4 & 8 & 473 & 7.74597 & 0.127017 & 0.37441344 & 790.3 \\
4.5 & 9 & 809 & 8.7178 & 0.141101 & 0.39978704 & 1,214.6 \\
5 & 10 & 1545 & 10 & 0 & 0.50093193 & 1,539.3 \\
5.5 & 11 & 2233 & 10.7703 & 0.114835 & 0.49450219 & 2,282.7 \\
6 & 12 & 3457 & 11.8322 & 0.0839202 & 0.54053601 & 2,938.5 \\
6.5 & 13 & 5001 & 12.8062 & 0.0968758 & 0.56771933 & 3,807.9 \\
7 & 14 & 7281 & 14 & 0 & 0.61451024 & 4,567.5 \\
7.5 & 15 & 9489 & 14.8324 & 0.0838015 & 0.60772296 & 6,125.0 \\
8 & 16 & 12833 & 15.8745 & 0.0627461 & 0.63489001 & 7,379.9 \\
8.5 & 17 & 16657 & 16.8523 & 0.0738502 & 0.64662328 & 9,103.0 \\
9 & 18 & 22409 & 18 & 0 & 0.69212206 & 9,968.2 \\
9.5 & 19 & 27233 & 18.868 & 0.0660189 & 0.67753435 & 12,961.3 \\
10 & 20 & 34569 & 19.8997 & 0.0501256 & 0.70051440 & 14,779.0 \\
100 & 200 & 476927289 & 199.99 & 0.005 & 0.96645675 & 16,552,931.1 \\
1000 & 2000 & 4918072576937 & 2000 & 0 & 0.99660987 &
16,729,623,607.7 \\
10000 & 20000 & 49331268889291561 & 20000 & 0 & 0.99966051 &
16,753,116,155,232.0 \\
\end{longtable}
}

The filling ratio shows local increases and decreases, while for large
\(R\) it tends toward \(1.0\). On the other hand, the gap corresponds to
the boundary layer of the 4-dimensional hyperball. The leading scale of
the boundary-layer volume is proportional to the surface area of the
4-dimensional ball. Therefore the leading order of the gap grows as
\(R^3\). The asymptotic constant and the detailed lattice oscillations
are left for separate analysis.

\begin{center}\rule{0.5\linewidth}{0.5pt}\end{center}

\subsection{6. Basic Check Values}\label{basic-check-values}

Extracting from Table 1 the first values at integer radius,

\[
N_0(1)=1,
\]

\[
N_0(2)=9,
\]

\[
N_0(3)=137,
\]

\[
N_0(4)=473,
\]

\[
N_0(5)=1545 .
\]

In particular,

\[
N_0(3)=137
\]

is obtained as the number of unit 4-cells of side 1 fully inscribed in
the 4-dimensional hyperball of radius 3.

This 137 is not introduced under any assumed correspondence with a
physical constant; it is a pure counting result obtained from the
4-dimensional integer lattice, radius 3, the unit cell, and the
full-inscription condition.

\begin{center}\rule{0.5\linewidth}{0.5pt}\end{center}

\subsection{\texorpdfstring{7. Decomposition at
\(R=3\)}{7. Decomposition at R=3}}\label{decomposition-at-r3}

For \(R=3\), the condition is

\[
\sum_{i=1}^{4}
\left(|k_i|+\frac12\right)^2
\le 9 .
\]

Multiplying both sides by 4,

\[
\sum_{i=1}^{4}
(2|k_i|+1)^2
\le 36 .
\]

The possible shells of sums of odd squares in this range are

\[
4,\;12,\;20,\;28,\;36 .
\]

The counts per shell are

\[
4:1,\quad
12:8,\quad
20:24,\quad
28:40,\quad
36:64 .
\]

Hence

\[
1+8+24+40+64=137 .
\]

Note that the shell \(28\) is the combined count of the \((3,3,3,1)\)
type and the \((5,1,1,1)\) type.

\begin{center}\rule{0.5\linewidth}{0.5pt}\end{center}

\subsection{8. Reproduction Algorithm}\label{reproduction-algorithm}

Table 1 can be reproduced by the following pseudocode.

\begin{Shaded}
\begin{Highlighting}[]
\KeywordTok{def}\NormalTok{ count\_cells(R):}
\NormalTok{    count }\OperatorTok{=} \DecValTok{0}
\NormalTok{    max\_s }\OperatorTok{=} \DecValTok{0}
\NormalTok{    max\_abs }\OperatorTok{=}\NormalTok{ floor(R }\OperatorTok{{-}} \FloatTok{0.5}\NormalTok{) }\OperatorTok{+} \DecValTok{1}

    \ControlFlowTok{for}\NormalTok{ k1 }\KeywordTok{in} \BuiltInTok{range}\NormalTok{(}\OperatorTok{{-}}\NormalTok{max\_abs, max\_abs }\OperatorTok{+} \DecValTok{1}\NormalTok{):}
      \ControlFlowTok{for}\NormalTok{ k2 }\KeywordTok{in} \BuiltInTok{range}\NormalTok{(}\OperatorTok{{-}}\NormalTok{max\_abs, max\_abs }\OperatorTok{+} \DecValTok{1}\NormalTok{):}
        \ControlFlowTok{for}\NormalTok{ k3 }\KeywordTok{in} \BuiltInTok{range}\NormalTok{(}\OperatorTok{{-}}\NormalTok{max\_abs, max\_abs }\OperatorTok{+} \DecValTok{1}\NormalTok{):}
          \ControlFlowTok{for}\NormalTok{ k4 }\KeywordTok{in} \BuiltInTok{range}\NormalTok{(}\OperatorTok{{-}}\NormalTok{max\_abs, max\_abs }\OperatorTok{+} \DecValTok{1}\NormalTok{):}
\NormalTok{            s }\OperatorTok{=} \BuiltInTok{sum}\NormalTok{((}\BuiltInTok{abs}\NormalTok{(ki) }\OperatorTok{+} \FloatTok{0.5}\NormalTok{)}\OperatorTok{**}\DecValTok{2} \ControlFlowTok{for}\NormalTok{ ki }\KeywordTok{in}\NormalTok{ [k1,k2,k3,k4])}
            \ControlFlowTok{if}\NormalTok{ s }\OperatorTok{\textless{}=}\NormalTok{ R}\OperatorTok{**}\DecValTok{2}\NormalTok{:}
\NormalTok{                count }\OperatorTok{+=} \DecValTok{1}
\NormalTok{                max\_s }\OperatorTok{=} \BuiltInTok{max}\NormalTok{(max\_s, s)}

\NormalTok{    rho }\OperatorTok{=}\NormalTok{ sqrt(max\_s) }\ControlFlowTok{if}\NormalTok{ count }\OperatorTok{\textgreater{}} \DecValTok{0} \ControlFlowTok{else} \DecValTok{0}
\NormalTok{    diagonal }\OperatorTok{=} \DecValTok{2} \OperatorTok{*}\NormalTok{ rho}
    \ControlFlowTok{return}\NormalTok{ count, rho, diagonal}
\end{Highlighting}
\end{Shaded}

This computation merely judges, for each cell, whether its farthest
vertex lies within the ball of radius \(R\).

\begin{center}\rule{0.5\linewidth}{0.5pt}\end{center}

\subsection{9. On the Absence of
Uncertainty}\label{on-the-absence-of-uncertainty}

This paper does not treat the weighted contribution of near-boundary
cells or the effective counting including the uncertainty width
\(\delta\).

That is, what is computed in this paper is the fully-inscribed counting
at

\[
\delta=0 .
\]

To perform counting with uncertainty, the indicator function must be
replaced by a weight function

\[
W_\delta(x) .
\]

However, how to determine the value of \(\delta\) or the form of that
function is undefined and outside the scope of this paper.

\begin{center}\rule{0.5\linewidth}{0.5pt}\end{center}

\subsection{10. Conclusion}\label{conclusion}

In this paper we computed, from radius \(0.5\) to \(10.0\) in steps of
\(0.5\), the number of unit 4-cells of side 1 on the 4-dimensional
integer lattice that are fully inscribed in the 4-dimensional hyperball
of radius \(R\).

The fully-inscribed cell count is defined by

\[
N_0(R)
=
\#\left\{
k\in\mathbb{Z}^4
\mid
\sum_{i=1}^{4}
\left(|k_i|+\frac12\right)^2
\le R^2
\right\} .
\]

This counting is purely geometric / integer-lattice in nature and
asserts no correspondence with physical constants.

In particular,

\[
N_0(1)=1,\qquad
N_0(2)=9,\qquad
N_0(3)=137
\]

are obtained.

The results of this paper provide a basic reference showing that the
sum-of-squares constraint, when the wavelength space or the frequency
space is chosen as a 4-dimensional lattice, can be treated as a concrete
unit-cell counting problem.

\begin{center}\rule{0.5\linewidth}{0.5pt}\end{center}

\subsection{Appendix A: CSV File}\label{appendix-a-csv-file}

Table 1 of this paper is also saved as the accompanying file

\texttt{paper2\_radius\_sweep\_R\_0\_5\_to\_10\_step\_0\_5.csv}.

The meaning of the columns is as follows.

\begin{itemize}
\tightlist
\item
  \texttt{R}: radius of the circumscribing ball
\item
  \texttt{sphere\_diameter\_2R}: diameter \(2R\) of the circumscribing
  ball
\item
  \texttt{full\_inclusion\_cell\_count\_N0}: fully-inscribed cell count
  \(N_0(R)\)
\item
  \texttt{stack\_vertex\_radius}: maximal vertex radius \(\rho(R)\) of
  the stacked cell set
\item
  \texttt{stack\_diagonal\_length}: stacked diagonal length \(2\rho(R)\)
\item
  \texttt{unused\_margin\_R\_minus\_stack\_radius}: radial margin
  \(R-\rho(R)\)
\item
  \texttt{shells\_S\_counts}: count per shell \(S=\sum_i(2|k_i|+1)^2\)
\end{itemize}

\begin{center}\rule{0.5\linewidth}{0.5pt}\end{center}

\subsection{Appendix B: Filling-Ratio and Gap Calculation
Program}\label{appendix-b-filling-ratio-and-gap-calculation-program}

The values of \(N_0(R)\), \(2\rho(R)\), the filling ratio \(V_1/V_0\),
and the gap \(V_0-V_1\) in Table 1 can be reproduced exactly by the
accompanying Python script

\texttt{paper2\_count\_fill\_gap.py}.

The naive quadruple loop in Section 8 has computational complexity
\(O(R^4)\) and becomes impractical for \(R=100\) and above. Therefore
the script rewrites the full-inscription condition as the integerized
positive odd-square condition

\[
\sum_{i=1}^{4}(2|k_i|+1)^2\le(2R)^2,
\]

constructs the one-axis odd-square distribution, convolves it into a
two-axis distribution \(g_2\), and then evaluates \(N_0(R)\) by
cumulative sums in roughly \(O(R^2)\). The value of \(2\rho(R)\) is
obtained by a two-pointer search over the keys of \(g_2\). All
computations are performed using integer arithmetic; the values up to
\(R=10000\) are exact rather than approximate. As a checksum, the script
internally confirms \(N_0(1{:}5)=1,9,137,473,1545\).

\end{document}
